\begin{apartats}

\apartat[1,25]{0,75}
Calcula $\begin{vmatrix}a-2&3\\4&a+2\end{vmatrix}$ i troba els valors de $a$ per als quals és nul.

\begin{solucio}
$(a-2)(a+2)-12=a^2-4-12=a^2-16$. És nul per a $a=4$ i per a $a=-4$.
\end{solucio}

\begin{tria}{determinant-lletres}
\itemtria{ordre-3-lletres}{1}{1,25}
Calcula el determinant $\begin{vmatrix}x&1&1\\1&x&1\\1&1&x\end{vmatrix}$, comprova que és igual a
$(x-1)^2(x+2)$ i digues per a quins valors de $x$ s'anul·la.

\begin{solucio}
Per Sarrus: $x^3+1+1-x-x-x=x^3-3x+2$. D'altra banda, $(x-1)^2(x+2)=(x^2-2x+1)(x+2)=x^3-3x+2$: coincideixen.
S'anul·la per a $x=1$ i per a $x=-2$.
\end{solucio}

\itemtria{familia}{1}{1,25}
Troba totes les matrius de la forma $\begin{pmatrix}a&b\\b&a\end{pmatrix}$ que tenen determinant nul, i
totes les que tenen determinant 1 amb $b=0$.

\begin{solucio}
$\begin{vmatrix}a&b\\b&a\end{vmatrix}=a^2-b^2=(a-b)(a+b)$. És nul si i només si $b=a$ o $b=-a$: les
matrius $\begin{pmatrix}a&a\\a&a\end{pmatrix}$ i $\begin{pmatrix}a&-a\\-a&a\end{pmatrix}$.\\
Amb $b=0$, el determinant és $a^2$, i val 1 si $a=1$ o $a=-1$: les matrius $I$ i $-I$.
\end{solucio}
\end{tria}

\begin{nomesllarg}
\apartat{0,75}
Calcula el determinant $\begin{vmatrix}2&-1&3\\1&4&-2\\0&5&1\end{vmatrix}$.

\begin{solucio}
Per Sarrus: $8+0+15-0+1+20=44$.
\end{solucio}
\end{nomesllarg}

\end{apartats}
